% !TEX root = ../../../../main.tex
% SECTION 3 --- Method.  Budget: 2.5 pages.
%
% BINDING RULES (rationale and history are in git):
%  1. Never read a fitted E as an entropy.  The NL floor moves 1.69 -> 1.82 on a
%     refit and the code floor 0.22 -> ~0 on a change of form, on the same runs.
%     Lean on the ORDERING of floors across channels, never on a value.
%  2. luo2025codescalingdatahungry is the code-vs-NL fit; yang2026codescaling is
%     WITHIN-code only and must never be cited for code-vs-English.
%  3. The exponent half of sec:method:scaling is a TRANSFER from nonparametric
%     estimation, not an imported theorem.  Keep the seam paragraph: the floor
%     claim is exact, the exponent claim is registered and separately
%     falsifiable.  An edit that blurs the two breaks the paper.
%  4. Do not restore the Shannon-band cap on KL(NL||pi_NL).  Falsified three
%     ways (reviews/2026-08-11-kl-nl-estimation-memo.md); every failure
%     flattered the paper.  It gets one footnote.

\section{Method}
\label{sec:method}

\subsection{One theorem, three renderings}
\label{sec:method:object}

In a dependently typed proof assistant a theorem is a closed term whose type is
the statement, so proving \(t\) and producing a term of type \(t\) are the same
act \citep{demoura2021lean4, barendregt2001proofassistants}.
We train on the serialisation \(\sigma\) of that term.
\(\sigma\) is \emph{injective}, so a distribution over strings is a distribution
over terms with no realisation freedom in between, and every quantity we compute
on the string channel is a quantity about terms.
Acceptance is decided rather than scored:
the kernel either accepts a candidate as a proof of \(t\) or does not, so
membership in the accept set \(A(t)\) is a fact about the candidate rather than
a judgement about it.
The soundness record --- one kernel-ratified verifier with a private certificate
constructor, twenty-two adversarial probes, and a retained upstream kernel
defect kept as a positive control --- is \cref{app:soundness}.

The same theorem admits several renderings, and this paper is about the
difference between them.
We use three:
the \textbf{language} rendering, an English statement and worked solution;
the \textbf{surface} rendering, the proof-assistant source a human types;
and the \textbf{term} rendering, the serialised kernel term the checker
consumes.
Autoformalisation is the map from the first to the last.

\subsection{Signal, noise, and where a prover's loss goes}
\label{sec:method:identity}

Fix a theorem set.
Each row carries a statement and a proof, and only the proof is ever scored.
Write \(Z\) for what the prover is handed --- the English statement together
with the formal type \(t\) of \cref{sec:method:object}, given to every arm alike
(R3, \cref{sec:method:protocol}) --- and NL and FL for the corpora of English
worked solutions and of the serialised kernel terms that certify them.
\(Z\) is the question;
the pair \((\mathrm{NL}, \mathrm{FL})\) is the answer written twice.
All quantities are summed nats per theorem, conditioned on \(Z\), and computed
on held-out rows, because a model scored on its training rows can memorise below
the entropy and void what follows silently (R7, \cref{sec:method:protocol}).
The conditioning on \(Z\) is suppressed in the displays below and restored in
\cref{eq:fibre}, where it does work.
We never work per token:
per-token quantities are properties of the tokenizer.

On each channel the held-out cross-entropy is corpus entropy plus model misfit
\citep{cover2006elements}:
\begin{equation}
\label{eq:cehkl}
  \CE_{\mathrm{NL}} = H(\mathrm{NL}) + \KL{\mathrm{NL}}{\pi_{\mathrm{NL}}},
  \qquad
  \CE_{\mathrm{FL}} = H(\mathrm{FL}) + \KL{\mathrm{FL}}{\pi_{\mathrm{FL}}}.
\end{equation}
The chain rule splits the English entropy against the formal side,
\(H(\mathrm{NL}) = I(\mathrm{NL};\mathrm{FL}) + H(\mathrm{NL}\mid\mathrm{FL})\),
with no premise consumed;
substituting into \cref{eq:cehkl} gives the decomposition the paper is built on:
\begin{equation}
\label{eq:identity}
  \CE_{\mathrm{NL}}
  \;=\; \underbrace{I(\mathrm{NL};\mathrm{FL})}_{\text{signal}}
  \;+\; \underbrace{H(\mathrm{NL}\mid\mathrm{FL})}_{\text{noise}}
  \;+\; \underbrace{\KL{\mathrm{NL}}{\pi_{\mathrm{NL}}}}_{\text{misfit}}.
\end{equation}
\Cref{fig:identity} is the whole accounting on one worked instance;
\cref{app:info} defines the quantities for readers who want them.

% FIGURE.  The picture lives in figures/identity_idiagram/body.tex and is
% shared, byte for byte, with the internal deck
% (publications/slides/2026-08-17/sections/figure.tex).  Design rationale and
% measured geometry travel in that component's README.
% CONTRACT for any other host: figures/identity_idiagram/requires.tex, plus
% \CE from flmnotation (which flmcore and flmbeamer both pull in).  The
% contract used to say \input{macros}; macros.tex was retired into
% flmnotation.sty on 2026-08-16.
\begin{figure}[t]
\centering
\input{figures/identity_idiagram/body}
\caption{The decomposition, on one worked instance.
Left: the same theorem twice, as English prose and as the tree of \Expr{}
constructors the kernel checks (\cref{app:expr});
the serialisation $\sigma$ reads the tree top-down and is injective, so a
distribution over streams is a distribution over terms.
Each bar is one channel's held-out cross-entropy as a length in summed nats per
theorem, and the bars are aligned \emph{on} their shared cell, because that cell
is the same quantity in both.
Each bar read left to right is noise, then signal, then misfit --- the top bar
is \cref{eq:identity}, whose terms we write signal-first;
the nested braces are \cref{eq:cehkl}; the shaded column read against
the bottom bar is \cref{eq:mathcap}.
\emph{Widths are schematic:} at the de-risk orders the shaded column would be
under one per cent of the top bar.}
\label{fig:identity}
\end{figure}

\paragraph{One definition, stated as one.}
We \emph{define} the mathematical content of a write-up as the information it
shares with the term that certifies it, \(I(\mathrm{NL};\mathrm{FL})\), and the
\emph{noise} of a rendering as what the certifying term leaves free,
\(H(\cdot\mid\mathrm{FL})\).
The motivation is propositions-as-types --- the kernel term \emph{is} a complete
derivation \citep{curry1934functionality, Howard1980} --- together with the
observation that for a prover the term is the deliverable.
Curry--Howard motivates the definition; nothing below rests on it as a premise.
The concession we flag rather than smuggle:
a kernel term may not capture everything a reader would call mathematics ---
motivation, intuition, choice of exposition --- so our ``noise'' contains, by
construction, whatever of that the term does not determine.
That is the honest attack on this framing, and the definition is our answer.

\paragraph{The headline inequality.}
Mutual information is at most either marginal entropy, and entropy is at most
cross-entropy, so
\begin{equation}
\label{eq:mathcap}
  I(\mathrm{NL};\mathrm{FL}) \;\le\; H(\mathrm{FL}) \;\le\; \CE_{\mathrm{FL}},
\end{equation}
and \(\CE_{\mathrm{FL}}\) is a number our harness reports per row:
a median of \(3.10\) nats per theorem in the de-risk pilot
(\cref{sec:experiments:pilot}), against an English objective of order
\(10^{2}\)--\(10^{4}\) nats \citep{shannon1951prediction,
deletang2024compression}.
At most \(\CE_{\mathrm{FL}}\) nats of a prover's loss can be signal;
everything above that line is noise or misfit.
Three properties do the work.
The cap assumes no floor on \(H(\mathrm{NL})\), so nothing is imported from the
entropy-of-English literature.%
\footnote{An earlier draft capped \(\KL{\mathrm{NL}}{\pi_{\mathrm{NL}}}\) using
Shannon's human-prediction band.
Withdrawn: a cap on the misfit is exactly an assumed floor on the entropy of
held-out mathematical prose, Shannon's \(0.6\) bits per character bounds
\(F_{100}\) rather than the entropy rate and presumes an ideal ranker
\citep{shannon1951prediction, coverking1978gambling}, and a measured \(0.500\)
bits per character on mathematical prose \citep{huang2024compression} already
sits below the band.
Every failure flattered the paper.}
The cap is ungameable in nats, since \cref{eq:mathcap} contains no
\(\pi_{\mathrm{NL}}\);
the \emph{fraction} \(\CE_{\mathrm{FL}}/\CE_{\mathrm{NL}}\) is gameable upward
by a bad English scorer, so we lead with nats and report the fraction across a
ladder of scorers.
And over-counting is safe:
anything inflating \(I\) or \(H(\mathrm{FL})\) --- shared surface correlates,
both sides running longer on harder problems, proof multiplicity --- loosens the
cap rather than invalidating it.

\paragraph{Why the operation is denoising.}
Under this definition the three renderings sit in a strict order.
The term channel leaves nothing free, \(H(\mathrm{FL}\mid\mathrm{FL}) = 0\);
the surface channel leaves some, because many scripts elaborate to one term;
the language channel leaves whatever English leaves free once the term is fixed,
which is the wording, the ordering, and the decision about which steps are worth
spelling out.
Autoformalisation walks that order end to end.
It is not a change of notation:
it removes the component of the description that carries no information about
the certifying term.
That operation is denoising, and the paper's question is how much of the
description it removes.

\paragraph{The one noise term we can enumerate.}
Read against the coarser object --- the theorem rather than a particular term ---
the term channel is not noiseless either:
several accepted terms realise one goal, and which appears is free.
This residual is the only noise term in the paper that is \emph{enumerable},
because \(A(t)\) is decidable, so we can sample the accept set and price the
freedom directly (\cref{sec:experiments:cap}).
No language channel permits that.
It also enters every bound above with a known sign, inflating
\(H(\mathrm{FL})\) and loosening \cref{eq:mathcap}, so the accounting never
needs it measured.

\paragraph{The absolute reading.}
The same chain rule bounds the noise from below,
\(H(\mathrm{NL}\mid\mathrm{FL}) = H(\mathrm{NL}) - I(\mathrm{NL};\mathrm{FL})
\ge H(\mathrm{NL}) - \CE_{\mathrm{FL}}\),
with no model on either side, since cross-entropy minus its own misfit is
identically \(H(\mathrm{NL})\).
Only the absolute statement --- that the noise is
\(\gtrsim \CE_{\mathrm{NL}} - \CE_{\mathrm{FL}}\) --- substitutes the measured
\(\CE_{\mathrm{NL}}\) for \(H(\mathrm{NL})\) and so needs the English misfit
small, and we defend that by measurement rather than citation:
cross-entropy \emph{differences} identify KL differences exactly, so a ladder of
scorers certifies \(\KL{\mathrm{NL}}{\pi_{\mathrm{weak}}} \ge
\CE(\pi_{\mathrm{weak}}) - \CE(\pi_{\mathrm{best}})\) with no assumption.
The residual bias has the friendly sign:
\(\CE_{\mathrm{NL}} \ge H(\mathrm{NL})\) always, so every improvement in the
English scorer makes the claim harder, not easier.

\subsection{A theory of the loss: renderings as fibres}
\label{sec:method:scaling}

The accounting above prices one corpus with one model.
A scaling law asks what \emph{any} model must pay.
This subsection is a mechanism that produces both, and it is what makes the two
instruments measurements of one thing rather than two adjacent results.
Write \(L_c(N,D)\) for the held-out loss of a model with \(N\) total
parameters trained on \(D\) tokens of channel \(c\), and fit the usual form \phy{Empirically this loss works for per-theorem and per nats. Here, I mean per-theorem.}
\citep{kaplan2020scaling, hoffmann2022chinchilla}
\begin{equation}
\label{eq:chinchilla}
  L_c(N,D) \;=\; E_c \;+\; A_c N^{-\alpha_c} \;+\; B_c D^{-\beta_c}.
\end{equation}

\paragraph{The construction.}
Write \(X_c\) for the string channel \(c\) presents, so \(X_{\mathrm{NL}}\) and
\(X_{\mathrm{FL}}\) are the two corpora of \cref{sec:method:identity} and
\(X_{\mathrm{surface}}\) is the third rung.
A rendering maps mathematical content to \(X_c\) and is not injective:
the content leaves choices open, and something has to make them.
Write \(\nu_c\) for the channel's \emph{realisation variable} --- everything the
content does not fix --- and read the certifying term \(X_{\mathrm{FL}}\) as the
content itself, which \(\sigma\) licenses, so that
\begin{equation}
\label{eq:fibre}
  X_c \;=\; r_c(X_{\mathrm{FL}}, \nu_c),
  \qquad
  H(X_c \mid Z) \;=\; \underbrace{I(X_c; X_{\mathrm{FL}} \mid Z)}_{\text{base}}
  \;+\; \underbrace{H(\nu_c \mid X_{\mathrm{FL}}, Z)}_{\text{fibre}}.
\end{equation}
The rendering map \(r_c\) is deterministic given content and realisation
variable, which is what turns the chain rule's residual into an entropy of
\(\nu_c\).
This is \cref{eq:identity} again, read structurally:
the signal is the base of the rendering and the noise is its fibre.
The three channels are three points on one tower, ordered by fibre.
The middle rung is not an assumption but a fact about the substrate:
elaboration is a genuine many-to-one function from surface source to kernel
term, so \(H(\nu_{\mathrm{surface}}\mid X_{\mathrm{FL}}) > 0\) is the entropy of
elaboration choices, and it is zero only for the term itself.
The English rung is the same construction with a fibre that is both far larger
and, unlike elaboration, not the fibre of a function --- a prose proof need not
determine a term at all.
Autoformalisation is the walk down this tower, and each step deletes a fibre.

\begin{proposition}[The floor is base plus fibre]
\label{prop:floor}
For every model, every \(N\) and every \(D\),
\(L_c(N,D) = H(X_c\mid Z) + \KL{X_c}{\pi} \ge H(X_c \mid Z)\).
The estimand of \(E_c\) is that limit, so by \cref{eq:fibre}
\(E_c\) is the channel's base plus its fibre entropy, and nothing else.
\end{proposition}

On the term channel the fibre is empty, so \(E_{\mathrm{FL}}\) is \emph{all}
base.
On the language channel \cref{eq:mathcap} caps the base at
\(\CE_{\mathrm{FL}}\), a few nats per theorem, so a floor of hundreds of nats
per theorem is essentially all fibre.

\paragraph{The second property is dimension, and it sets the exponent.}
Fibre entropy is how much is free.
Fibre \emph{dimension} is how much structure that freedom has --- how many
coordinates a model must resolve before it can predict form well --- and it is a
different quantity doing a different job.
In the resolution-limited regime \citep{bahri2024explaining} a model with \(N\)
parameters resolves a target manifold of intrinsic dimension \(d\) at spacing
\(N^{-1/d}\), and a smooth target approximated at that spacing costs squared
error \(N^{-4/d}\), so the model exponent goes as \(\alpha \approx 4/d\) ---
verified by measuring \(d\) and \(\alpha\) independently
\citep{sharma2022manifold}, and derived in \cref{app:noise}.
A rendering's manifold is locally a product of a content factor and a form
factor, and the intrinsic dimension of a product is the sum, so
\(d_c = d_{\mathrm{math}} + d_{\mathrm{form}}(c)\) and
\begin{equation}
\label{eq:alphafibre}
  \frac{1}{\alpha_c}
  \;\approx\; \frac{d_{\mathrm{math}} + d_{\mathrm{form}}(c)}{\kappa},
  \qquad \kappa \text{ channel-independent } (\approx 4).
\end{equation}
\emph{Reciprocal model exponents are affine in the fibre dimension.}
That is a sharper claim than ``noise flattens the exponent'':
it says the three channels' reciprocal exponents should be ordered \emph{and}
that their gaps measure the dimensions each rendering adds, with the term
channel as the zero-fibre reference.

The discrete-skill route gives the same conclusion by a different mechanism.
If learning decomposes into quanta acquired in decreasing order of use
frequency, a power law in use frequency produces the power law in loss
\citep{michaud2023quantization}.
A rendering's quanta are the content quanta together with the form quanta ---
that this register says ``hence'', that this house style displays its
equations --- and the ranked spectrum of a union is governed at large rank by
whichever population decays more slowly, so
\(\alpha_c \approx \min(\alpha_{\mathrm{content}}, \alpha_{\mathrm{form}}(c))\).
Deleting the form population can only raise the exponent.
The classical rates say the same thing in the limit:
a smooth target costs \(n^{-2s/(2s+d)}\) from noisy observations against
\(n^{-2s/d}\) from exact ones \citep{stone1982optimalrates,
tsybakov2009nonparametric}, and under label noise the fast rate survives only to
a crossover \citep{cui2021noisycrossover}.
Three routes, one conclusion, and the same \(d\) in two of them.

\paragraph{Two properties, two dials, and they separate.}
A fibre can be high-entropy and low-dimension:
an independent coin flip between two synonyms raises \(E\) and leaves \(\alpha\)
alone.
It can equally be low-entropy and high-dimension:
a house style is nearly deterministic given enough context, so it costs little
entropy, but a model must resolve it, and that costs exponent.
So the floor and the exponent are consequences of \emph{different properties of
the same fibre}, and an experiment can move them independently --- generator
temperature adds fibre entropy, generator style diversity adds fibre dimension.
\Cref{sec:experiments:law} registers that as the theory's discriminating test,
because no account on which noise is a single scalar can pass it.

\paragraph{What the theory already has to explain.}
Fitted head to head under one functional form, code carries both a lower floor
than natural language (\(0.22\) against \(\num{scaling_fit/nl_ref/E}\) nats per token) \emph{and} a
steeper model exponent (\(0.485\) against \(\num{scaling_fit/nl_ref/alpha}\))
\citep{luo2025codescalingdatahungry, hoffmann2022chinchilla}.
Two observations, one mechanism:
a smaller fibre has less entropy and fewer dimensions at once.
Read through \cref{eq:alphafibre} the implied dimensions are \(8.2\) and
\(11.8\), so English carries roughly four content-irrelevant dimensions that
code does not --- a number of the right order for this literature, and one our
own three channels re-measure under matched content.
The same comparison also \emph{restricts} the theory, which is why we state it
here rather than in a limitations paragraph.
Luo's data exponents barely move (\(0.298\) against \(\num{scaling_fit/nl_ref/beta}\)) while the model
exponents move a great deal.
\Cref{eq:alphafibre} is a statement about \(\alpha\), the exponent
\citet{sharma2022manifold} derive;
a version of this theory predicting that \(\beta\) moves in step is already
refuted by those numbers, and we therefore register the \(\alpha\) prediction
and leave \(\beta\) open.
Within code, stricter languages carry lower floors across a thousand runs
\citep{yang2026codescaling};
the theory predicts they carry steeper model exponents too, which their released
per-language fits can settle without any new training, and we pre-register that
check.

\paragraph{Why the published fits cannot carry the mechanism.}
Both orderings run in the predicted direction, and neither is evidence for the
fibre account, because in both designs the corpora differ in \emph{content} as
well as in fibre.
There is no common object that a Python corpus and an English corpus are two
renderings of, so nothing separates ``smaller fibre'' from ``easier, more
repetitive, more templated'' --- \(d_{\mathrm{math}}\) is not held fixed.
Our design fixes it by construction:
the three channels carry the \emph{same theorems}, joined row by row, so
\(d_{\mathrm{math}}\) and the base term are common to all three arms and every
difference between fitted floors and exponents is a difference in fibre.
The published ordering is an association;
this is the controlled contrast, and it is why the experiment is worth running
even though its direction is already in the literature.

\paragraph{An apparent contradiction, resolved.}
The same comparison reports code as \emph{more} data-hungry at the
compute-optimal allocation, which reads as the opposite of ``smaller fibre, less
data''.
It is not.
Minimising \cref{eq:chinchilla} at fixed compute \(C \propto ND\) gives
\(D_{\mathrm{opt}}/N_{\mathrm{opt}} \propto
C^{(\alpha-\beta)/(\alpha+\beta)}\) (\cref{app:noise}), an exponent of
\(0.10\) for the natural-language fit and \(0.24\) for the code fit.
Code's optimal data-to-parameter ratio grows faster with compute \emph{because}
its model exponent is steeper --- the very fact \cref{eq:alphafibre} predicts.
The reported data-hunger is a budget-splitting statement and is downstream of
the theory rather than in tension with it.

\paragraph{Two disciplines, both binding.}
\emph{Units.}
We fit per token, because that is what \cref{eq:chinchilla} is for, then convert
each floor by that channel's tokens-per-theorem before comparing.
A floor in nats per token is a statement about a tokenizer as much as about a
corpus, and a channel that compresses its own vocabulary wins by construction.
If the predicted ordering exists only per token, we have measured our
tokenizers, and \cref{sec:experiments:law} treats that as a kill condition
rather than a caveat.
\emph{Estimator, not estimand.}
\Cref{prop:floor} concerns the limit of the achievable loss;
the fitted \(E_c\) is a known unstable estimator of it, moving
\(\num{scaling_fit/nl_ref/E} \to 1.82\) under a refit \citep{besiroglu2024chinchilla} and
\(0.22 \to \approx 0\) under a change of functional form on the same runs
\citep{luo2025codescalingdatahungry}, and the work that introduced the reading
warns against identifying the irreducible loss with an entropy
\citep{henighan2020scaling}.
We therefore fit all three channels under one form on one grid and read
\emph{orderings and gaps}, never a value.
\subsection{The matched protocol}
\label{sec:method:protocol}

The comparisons are only worth making if they are constructed to be fair, and
the fairness is a protocol rather than an argument.
We pre-register it in the form fixed before any training compute was spent.
\textbf{R1, one base model or none}:
the cap arms start from one recorded snapshot of the same pretrained base model;
the scaling arms are from scratch on every rung, because a pretrained body
imports an unknown quantity of the very data the fit is about.
\textbf{R2, one theorem set}:
same declarations, joined by problem identity, with the join rate and the full
cascade published.
\textbf{R3, identical conditioning}:
every arm scored as conditional generation given the \emph{byte-identical}
English statement, and no arm given the formal statement unless all are --- the
condition the predecessor's headline figure failed, which alone can account for
a large entropy gap.
\textbf{R4, one scorer, one unit}:
summed nats per theorem from one routine; per-token loss is a fitting device and
a diagnostic, never a cross-channel comparison.
\textbf{R5, comparable convergence, published}:
each arm trained to its own early-stopping criterion on a held-out slice of its
own channel, with all learning curves printed.
\textbf{R6, tokenizer replication}:
the term arm scored twice, under its native vocabulary and under the base
model's, both numbers printed;
we do not claim per-theorem nats are exactly tokenizer-invariant, only that the
comparison does not turn on it.
\textbf{R7, decontamination and provenance, declared first}:
the evaluation slice screened against the standard suites and near-duplicate
scanned against the public sources of its English side;
every row labelled human-authored or machine-generated, strata never pooled
silently;
and any fine-tuned scorer fine-tuned on train and scored on held-out rows only,
because memorisation drives a measured \(\CE\) below \(H\) and voids
\cref{eq:cehkl} silently.
One further admissibility rule --- that an accept may cite only constants already
elaborated when the goal's own declaration was --- is stated with its
measurements in \cref{app:admissibility}.
